\documentclass[aspectratio=169]{beamer}
\usetheme{floral}
\florasetup{layout=surround, seed=3, density=1.2}
\setbeamersize{text margin left=2.5cm, text margin right=1cm}

\title{beamer-floral}
\subtitle{Procedural watercolour margins}
\author{Documentation}
\institute{beamer-floral}
\date{}

\begin{document}

\florathispage{layout=left, density=2.0, mix={leaves=3, leaflets=3, sprays=3, bloom=3}, tune={gesture={top_inset=0}}}
\begin{frame}[plain]
  \titlepage
\end{frame}

\section{Palettes}
\subsection{Warm}

\begin{frame}
  \frametitle{Warm palette}
  \begin{block}{Gold hero, warm blooms}
    The warm palette draws a gold hero, warm blooms and earth foliage.
  \end{block}
  \begin{itemize}
    \item The text block is measured.
    \item The flora grows around it.
  \end{itemize}
  \begin{alertblock}{Alert}
    \alert{Attention} uses yama-budo.
  \end{alertblock}
  \begin{exampleblock}{Example}
    The example uses yama-guri.
  \end{exampleblock}
  \begin{description}
    \item[Label] Description labels use the block accent.
  \end{description}
\end{frame}

\florapalette{blue}
\subsection{Blue}

\begin{frame}
  \frametitle{Blue palette}
  \begin{block}{Azure blooms}
    The blue palette keeps blooms, leaves and leaflets in the iroshizuku blues.
  \end{block}
  \begin{itemize}
    \item Blue blooms and blue foliage share one ink family.
    \item The gold hero and sprays stay warm.
  \end{itemize}
  \begin{alertblock}{Alert}
    \alert{Attention} uses murasaki-shikibu.
  \end{alertblock}
  \begin{exampleblock}{Example}
    The example uses tsuki-yo.
  \end{exampleblock}
  \begin{description}
    \item[Label] Description labels use kon-peki.
  \end{description}
\end{frame}

\florapalette{green}
\subsection{Green}

\begin{frame}
  \frametitle{Green palette}
  \begin{block}{Blue blooms, green leaves}
    The green palette pairs blue blooms with green foliage.
  \end{block}
  \begin{itemize}
    \item Leaves and leaflets take shin-ryoku and chiku-rin.
    \item Alerts keep their own violet from edo-murasaki.
  \end{itemize}
  \begin{alertblock}{Alert}
    \alert{Attention} uses edo-murasaki.
  \end{alertblock}
  \begin{exampleblock}{Example}
    The example uses chiku-rin.
  \end{exampleblock}
  \begin{description}
    \item[Label] Description labels use shin-ryoku.
  \end{description}
\end{frame}

\florapalette{warm}
\section{Layout and density}
\subsection{Layout}

\florathispage{layout=left}
\begin{frame}
  \frametitle{Layout: left}
  \framesubtitle{Floral margin only}
  \begin{block}{A single margin}
    The left layout grows the border in the left margin and across the top and foot.
  \end{block}
  \begin{itemize}
    \item The text block stays clear.
    \item The right margin is paper.
  \end{itemize}
\end{frame}

\florathispage{layout=surround}
\begin{frame}
  \frametitle{Layout: surround}
  \framesubtitle{A border around the text}
  \begin{block}{A closed border}
    The surround layout draws a crown above the text and a foot below it.
  \end{block}
  \begin{itemize}
    \item The crown stops under the running head.
    \item The palette supplies the block and alert inks.
  \end{itemize}
\end{frame}

\subsection{Density}

\florathispage{density=0.6}
\begin{frame}
  \frametitle{Density: 0.6}
  \framesubtitle{Airier}
  \begin{block}{Less growth}
    A lower density leaves more paper between the stems and blooms.
  \end{block}
  \begin{itemize}
    \item Sprays and leaves thin out.
    \item The same seed keeps the shape.
  \end{itemize}
\end{frame}

\florathispage{density=2.0}
\begin{frame}
  \frametitle{Density: 2.0}
  \framesubtitle{Fuller}
  \begin{block}{More growth}
    A higher density fills the margin with more stems, leaves and sprays.
  \end{block}
  \begin{itemize}
    \item The layout stays the same.
    \item The text block stays clear.
  \end{itemize}
\end{frame}

\end{document}
